\begin{abstract}
Reproducing a known IoT vulnerability requires more than generating an
exploit: an analyst must obtain the affected firmware, establish a usable
execution environment, and connect the observed behavior to the specified
target. We investigate how LLM agents perform and behave across this
workflow. Our study uses an archive of 450 runs spanning 30 CVEs, five
model configurations, and three repetitions. We organize the analysis
around reproduction outcomes, execution obstacles, and agents' subsequent
responses, using target provenance, execution artifacts, and trajectories
rather than an aggregate capability score. Initial archive inspection
identifies candidate real-component executions alongside demonstrations
against reimplemented services, illustrating why executable output alone
does not establish reproduction. \todo{ABS-01}{Insert the audited outcome
distribution and two principal cross-case findings after evidence review.}
We distinguish component-level triggering from reproduction through the
documented entry point, and separate declared simulation from unsupported
success claims. The study aims to clarify which parts of autonomous
reproduction are observed, which remain untested after early failures,
and what evidence is needed to interpret an agent's report.
\end{abstract}
